% ============================================================
%  LNCS template — Truncated Differential + BCT Key Recovery
% ============================================================
\documentclass[runningheads]{llncs}

\usepackage[T1]{fontenc}
\usepackage{amsmath,amssymb,amsthm}
\usepackage{mathtools}          % \coloneqq, etc.
\usepackage{booktabs}           % professional tables
\usepackage{xcolor}
\usepackage{hyperref}
\hypersetup{
  colorlinks=true,
  linkcolor=black,
  citecolor=black,
  urlcolor=black
}
\usepackage{microtype}          % better typography
\usepackage{graphicx}
\usepackage{algorithm}
\usepackage{algpseudocode}

% ── Notation shorthands ──────────────────────────────────────
\newcommand{\Ftwon}{\mathbb{F}_2^n}
\newcommand{\Z}{\mathbb{Z}}
\newcommand{\DDT}{\mathrm{DDT}}
\newcommand{\BCT}{\mathrm{BCT}}
\newcommand{\tBCT}{\mathrm{BCT}_{\mathrm{trunc}}}
\newcommand{\pself}{p^{\mathrm{self}}_{\mathrm{trunc}}}
\newcommand{\WKP}{\mathrm{WKP}}
\newcommand{\LLR}{\mathrm{LLR}}
\newcommand{\speck}{\textsc{Speck}32/64}
\newcommand{\xor}{\oplus}
\newcommand{\ror}[2]{(#1 \ggg #2)}
\newcommand{\rol}[2]{(#1 \lll #2)}

% ── Theorem-like environments ────────────────────────────────
\newtheorem{definition}{Definition}
\newtheorem{theorem}{Theorem}
\newtheorem{proposition}{Proposition}
\newtheorem{lemma}{Lemma}
\newtheorem{corollary}{Corollary}
\newtheorem{remark}{Remark}

% ============================================================
\begin{document}
% ============================================================

\title{Theoretical Key Recovery from Truncated Differentials:\\
       A BCT-Algebraic Wrong-Key Profile for \speck{}}

% ── Suppress author details for anonymous review ─────────────
% Comment out the \author / \authorrunning / \institute block
% below and uncomment the anonymous version for submission.
\author{Anonymous Author(s)}
\authorrunning{Anonymous}
\institute{}

% ── De-anonymised version (uncomment for camera-ready) ───────
%\author{FirstName LastName\inst{1} \and
%        FirstName LastName\inst{2}}
%\authorrunning{LastName and LastName}
%\institute{Institution Name, City, Country
%           \email{email@institution.edu}
%  \and
%           Institution Name, City, Country
%           \email{email@institution.edu}}

\maketitle

% ============================================================
\begin{abstract}
Neural-network-based distinguishers have achieved state-of-the-art
results in differential cryptanalysis of \speck{}, most notably
through Gohr's CRYPTO~2019 attack framework.  That framework rests on
two black-box components—a deep residual distinguisher and an
empirically sampled wrong-key response profile—whose internal
structure has so far resisted closed-form analysis.  In this paper we
replace both components with algebraically defined counterparts and
thereby provide a fully interpretable reconstruction of the attack.
We substitute the neural distinguisher with a truncated differential
distinguisher whose scores are computed directly from the cipher's
difference distribution table (DDT), and we derive the wrong-key
profile analytically via a novel application of the boomerang
connectivity table (BCT).  Specifically, for each 16-bit subkey
offset~$g$ we show that the expected distinguisher score under the
corresponding wrong-key hypothesis equals a \emph{truncated
self-returning BCT probability} $\pself(g)$, which admits an exact
bit-by-bit computation through a carry-propagation transfer matrix.
The resulting profile is derived rather than measured, its shape is
governed by the algebraic structure of \speck{}'s modular-addition
round function, and it requires no neural network training or
large-scale sampling.  We demonstrate that the reconstructed Bayesian
beam-search framework achieves key-recovery performance competitive
with Gohr's original attack while being amenable to complexity
analysis.  Our results suggest that the statistical learning structure
underlying neural-distinguisher attacks on ARX ciphers can in many
cases be recovered analytically from classical differential tools.
\keywords{Speck32/64 \and differential cryptanalysis \and truncated
differentials \and boomerang connectivity table \and wrong-key profile
\and key recovery \and ARX cipher}
\end{abstract}

% ============================================================
\section{Introduction}
\label{sec:intro}
% ============================================================

Differential cryptanalysis, introduced by Biham and
Shamir~\cite{BS91}, exploits the non-uniform propagation of XOR
differences through a block cipher's round function to recover secret
key material.  For an $n$-bit block cipher, the fundamental object
governing this propagation is the \emph{difference distribution table}
(DDT): an array indexed by input and output differences whose entries
record the probability that a uniformly random input pair with the
given input difference produces the given output difference after one
round of encryption.  The boomerang attack, introduced by
Wagner~\cite{Wag99} and subsequently refined through the notion of the
\emph{boomerang connectivity table} (BCT) by Cid et
al.~\cite{CHPSS18}, extends the differential framework to cover
compatible transitions across the encryption and decryption
directions.  Both tables have become standard tools in the security
evaluation of symmetric-key primitives.

\subsubsection{Speck32/64 and classical attacks.}
The \speck{} cipher~\cite{BRRSSST13} belongs to the ARX (Add-Rotate-XOR)
design family and has attracted sustained cryptanalytic interest as a
lightweight benchmark target.  Its round function consists of a
modular addition, a rotation, and an XOR, combined with a key
injection:
\begin{equation}
  \label{eq:speck-round}
  (x,\,y) \;\mapsto\;
  \bigl(\ror{x}{7} \boxplus y \xor k,\;\;
        \rol{y}{2} \xor (\ror{x}{7} \boxplus y \xor k)\bigr),
\end{equation}
where $\boxplus$ denotes addition modulo $2^{16}$ and $k$ is the
current round subkey.  The differential behaviour of modular
addition is well understood~\cite{LM01}, enabling precise computation
of DDT and BCT entries for individual rounds of \speck{}.  Classical
differential attacks have been mounted on up to 8 rounds of \speck{}
using optimised distinguishers~\cite{ALLW14,KLT15}, but extending
these attacks further has proven difficult due to the rapid diffusion
of differences across rounds.

\subsubsection{Gohr's neural-network attack.}
In a landmark result at CRYPTO 2019, Gohr~\cite{Goh19} demonstrated
that deep residual networks could serve as highly effective
distinguishers for 7- and 8-round \speck{}, and embedded them within a
sophisticated Bayesian key-recovery framework to mount practical attacks
on up to 11 rounds.  The framework integrates three tightly coupled
components.

\begin{enumerate}
  \item \textbf{Neural distinguisher.}  A deep residual network is
    trained to classify ciphertext pairs produced under a fixed input
    difference from randomly generated pairs.  Given a ciphertext pair
    $(C, C')$ with $C \xor C' = \Delta_{\mathrm{out}}$, the network
    outputs a probability $p \in (0,1)$ that the pair arises from a
    real differential rather than from two independent random
    plaintexts.  The per-pair log-likelihood ratio (LLR)
    $\log_2(p/(1-p))$ is accumulated over all available ciphertext
    pairs to produce a score for each subkey hypothesis.

  \item \textbf{Wrong-key response profile (WKP).}  For each of the
    $2^{16}$ possible last-round subkey offsets $g = k_r \xor k_g$
    (where $k_r$ is the real subkey and $k_g$ is the guessed subkey),
    Gohr measures the \emph{mean} and \emph{standard deviation} of the
    network's output when ciphertext pairs are partially decrypted
    under the wrong subkey $k_g$.  This profiling step requires on the
    order of $10^6$ forward passes through the network for each of the
    $2^{16}$ offset values, making it computationally intensive and
    inherently tied to the specific trained network.

  \item \textbf{Bayesian beam search.}  Given a set of ciphertext
    pairs, Gohr iterates a beam search over the subkey space.  In each
    iteration, the current pool of $B = 32$ subkey candidates is
    partially decrypted and fed to the neural network; the resulting
    empirical output means are compared to the precomputed WKP via a
    Bayesian distance metric (an $\ell^2$ residual against the profile),
    and the $B$ candidates whose profile best matches the observed
    statistics are retained for the next iteration.  A two-layer
    nesting recovers the last two round subkeys sequentially, followed
    by a local refinement step that exhaustively explores candidates
    within Hamming distance 2 of the current best guess.
\end{enumerate}

This framework achieved results far beyond what had been attained with
classical methods and demonstrated, for the first time, that
machine-learning-based distinguishers could be integrated into a full
key-recovery attack with practical data and time complexity.

\subsubsection{The interpretability gap.}
Despite its empirical success, Gohr's framework exhibits a fundamental
limitation: it is constructed almost entirely from observations rather
than proofs.  The neural distinguisher is a black box whose internal
representations have no direct cryptanalytic interpretation.  The WKP
is produced by brute-force sampling and carries no closed-form
description.  As a consequence:
\begin{itemize}
  \item No theoretical bound governs the \emph{shape} of the WKP, so
    the signal-to-noise ratio of the beam search cannot be predicted
    analytically.
  \item No complexity formula characterises the minimum number of
    ciphertext pairs required for the search to succeed with a given
    probability.
  \item Generalising the attack to a related cipher, or to a different
    number of rounds, requires repeating the same expensive empirical
    construction—training a new network and resampling the
    profile—from scratch.
\end{itemize}
These limitations have prompted a line of follow-up work aimed at
understanding why neural distinguishers work as well as they
do~\cite{BM21,BGMR23,HWW23,LLLW22}, but a complete analytical
reconstruction of Gohr's full key-recovery framework has not yet been
achieved.

\subsubsection{This work: a BCT-algebraic reconstruction.}
We propose a theoretically grounded reconstruction of Gohr's key
recovery framework in which both black-box components are replaced by
algebraically defined counterparts derived from the differential
structure of \speck{}.

\paragraph{Truncated differential distinguisher.}
We replace the neural network with a \emph{truncated differential
distinguisher}.  Rather than learning patterns from training data, the
distinguisher assigns to each ciphertext pair a LLR computed directly
from the DDT of reduced-round \speck{}: after partially decrypting the
pair by one round under the candidate subkey, we evaluate the
log-probability of the resulting truncated output difference under the
target differential distribution.  The truncation—fixing only a subset
of the 32 output-difference bits—allows us to balance statistical
power against computation: a coarser truncation reduces the
distinguisher's information content but yields a simpler,
smaller-footprint probability table.  Every score is an explicit
function of the observed output difference and the precomputed
differential probabilities, with no black-box inference involved.

\paragraph{BCT-algebraic wrong-key profile.}
The central analytical contribution of this work is a closed-form
derivation of the WKP.  We show that, for a one-round
truncated differential distinguisher applied to \speck{}, the
expected LLR score under wrong-key hypothesis $g = k_r \xor k_g$ is
determined by the probability that the truncated differential
\emph{self-returns} under key offset~$g$:
\begin{equation}
  \label{eq:pself-informal}
  \pself(g) \;\coloneqq\;
  \Pr_x\!\bigl[\,
    \Delta_{\mathrm{out}} \in \mathcal{T}
    \;\big|\;
    \Delta_{\mathrm{in}} \in \mathcal{T},\;
    \text{one round of \speck{} with key offset } g
  \,\bigr],
\end{equation}
where $\mathcal{T}$ is the truncated class defined by the chosen mask.
We show that $\pself(g)$ can be computed \emph{exactly} for every
$g \in \{0,1\}^{16}$ via a BCT-style algebraic transfer-matrix
computation that processes each bit position independently and
propagates carry distributions through a small $16 \times 16$
stochastic matrix $M$.  The resulting profile requires no network
training and no sampling; it is a direct consequence of the algebraic
structure of the ARX round function.  In particular, the non-trivial
variation of $\pself(g)$ across key offsets—the phenomenon that makes
the Bayesian beam search effective—is explained entirely by the
differential properties of modular addition and the specific
bit-pattern of the chosen truncated differential.

\paragraph{Interpretable beam-search analysis.}
With both components in closed form, we provide a theoretical
characterisation of the Bayesian beam search.  The signal available to
the search is the gap $\pself(0) - \pself(g)$ between the correct-key
and wrong-key self-returning probabilities; its magnitude and
distribution over $g$ govern the convergence rate.  We derive a
lower bound on the expected score advantage of the correct subkey, and
use it to bound the data complexity of the attack.  We further show
that the UCB-guided structure selection adopted by Gohr can be
understood as a multi-armed bandit operating on noisy estimates of
these same self-returning probabilities.

\subsubsection{Contributions.}
We summarise the contributions of this paper as follows.
\begin{enumerate}
  \item We introduce the \emph{truncated self-returning BCT
    probability} $\pself(g)$ and show that it constitutes the
    wrong-key profile of any one-round truncated differential
    distinguisher applied to an ARX cipher whose round function
    contains a single modular addition (Sect.~\ref{sec:bct}).

  \item We present an exact bit-serial transfer-matrix algorithm for
    computing $\pself(g)$ for all $g \in \{0,1\}^{16}$ in time
    $O(16 \cdot 2^{16})$, requiring no sampling
    (Sect.~\ref{sec:algorithm}).

  \item We instantiate the framework on \speck{} with two carefully
    selected truncated differentials (with input differences
    $\mathtt{0x181E}$ and $\mathtt{0xC6A5}$) and verify that the
    computed profiles match empirical measurements to within
    statistical noise (Sect.~\ref{sec:experiments}).

  \item We integrate the truncated distinguisher and the algebraic WKP
    into a full reconstruction of Gohr's Bayesian key-recovery attack
    and demonstrate key recovery competitive with the original
    neural-network-based framework on 9--11 rounds of \speck{}
    (Sect.~\ref{sec:attack}).

  \item We derive data-complexity bounds for the beam search in terms
    of $\pself$ and compare them to Gohr's empirically observed data
    usage (Sect.~\ref{sec:complexity}).
\end{enumerate}

\subsubsection{Related work.}
Neural distinguishers for block ciphers have been studied extensively
since Gohr's breakthrough.  Baksi et al.~\cite{BM21} explored similar
approaches for other ciphers; Benamira et al.~\cite{BGMR23} gave a
partial theoretical explanation of why neural networks can distinguish
$r$-round \speck{} by identifying the specific statistical features
learned by the network; Hou et al.~\cite{HWW23} improved distinguisher
accuracy through enhanced training procedures; and Liu et
al.~\cite{LLLW22} studied the connection between neural distinguishers
and classical DDT-based distinguishers.  Closest to our work is the
independent observation by several authors that the neural distinguisher
at 5 and 6 rounds is closely approximated by a DDT-based classifier;
however, no prior work has derived the wrong-key profile algebraically
or connected it to the BCT.  On the BCT side, the original
formulation~\cite{CHPSS18} and its extensions~\cite{SBBB19,WQSN19}
focus on the probability of boomerang quartets in the middle rounds of
a cipher; our use of BCT-style algebra to compute single-round
wrong-key profiles is, to the best of our knowledge, new.  Truncated
differentials were introduced by Knudsen~\cite{Knu94} and have been
applied to \speck{} by several works~\cite{ALLW14,KLT15}; here we
combine them with the Bayesian key-ranking apparatus of Gohr's
framework for the first time.

\subsubsection{Organisation.}
Section~\ref{sec:prelim} reviews \speck{}, the DDT, the BCT, and
Gohr's original framework in the notation used throughout the paper.
Section~\ref{sec:bct} introduces the truncated self-returning BCT
probability and proves its role as the wrong-key profile.
Section~\ref{sec:algorithm} presents the transfer-matrix algorithm for
computing $\pself$.
Section~\ref{sec:attack} describes the full reconstructed key-recovery
attack.
Section~\ref{sec:complexity} derives data-complexity bounds.
Section~\ref{sec:experiments} presents experimental results.
Section~\ref{sec:conclusion} concludes and discusses open problems.

% ============================================================
%  Placeholder sections (to be expanded)
% ============================================================

\section{Preliminaries}
\label{sec:prelim}
[To be written.]

\section{Truncated Self-Returning BCT Probability}
\label{sec:bct}
[To be written.]

\section{Transfer-Matrix Algorithm}
\label{sec:algorithm}
[To be written.]

\section{Reconstructed Key-Recovery Attack}
\label{sec:attack}
[To be written.]

\section{Data Complexity Analysis}
\label{sec:complexity}
[To be written.]

\section{Experiments}
\label{sec:experiments}
[To be written.]

\section{Conclusion}
\label{sec:conclusion}
[To be written.]

% ============================================================
\bibliographystyle{splncs04}
\bibliography{refs}

\end{document}
